\documentclass[10pt]{beamer}
\usetheme{metropolis}

\usepackage[utf8]{inputenc}
\usepackage[russian,french]{babel}
\usepackage[T1]{fontenc}
\usepackage{geometry}
\geometry{
    left=2.5mm,
    right=2.5mm,
    top=0mm,
    bottom=0mm
}
\usepackage{amssymb} % pour \blackdiamond
\usepackage{tikz}
\usepackage{graphicx,animate}
\usepackage[dvipsnames]{xcolor}
\usepackage{listings}
\usepackage[table,xcdraw]{xcolor}
\usepackage{minted}
\usepackage{hyperref}

\lstset{%
basicstyle=\small\ttfamily,%
columns=flexible,%
breaklines=true,%
language=xml,%
numberstyle=\footnotesize,%
numbers=left,%
tabsize=2,%
title=\lstname,%
escapeinside={\%*}{*)},%
breaklines=true,%
breakatwhitespace=true,%
% framextopmargin=1em,%
% framexbottommargin=1em,%
moredelim=[s][\bfseries\color{blue}]{<}{\ },%
moredelim=[s][\bfseries\color{blue}]{</}{>},%
morecomment=[s]{<?}{?>},%
morecomment=[s]{<!--}{-->},%
commentstyle=\color{ForestGreen},%
stringstyle=\color{blue},%
extendedchars=false,%
inputencoding=utf8%
}

\newcommand{\replacementchar}{%
  \tikz[baseline=-0.6ex]{
    %\node[inner sep=0pt] (d) {\Large$\blackdiamond$};
    \node[inner sep=0pt] (d) {\Large$\blacksquare$};
    \node[white, font=\scriptsize\bfseries] at (d.center) {?};
  }%
}
\usepackage{tabularx}
%\usepackage[style=alphabetic,citestyle=authoryear]{biblatex}
%\addbibresource{bibliographie.bib}

\usepackage{xspace}
\newcommand{\themename}{\textbf{\textsc{metropolis}}\xspace}
\title{M1SOL041: Représentation et traitement des documents électroniques\\
Cours 2: Introduction au XML et au XML-TEI}

\date{}
\author{Marine Tiger\\  \quad {marine.tiger@sorbonne-universite.fr\\}}

\institute{Littératures françaises et comparée ED19, CELLF, Sorbonne Université}

\begin{document}
\maketitle
\setbeamertemplate{section in toc}[sections numbered]


\begin{frame}{Déroulé des séances}

% Please add the following required packages to your document preamble:
% \usepackage[table,xcdraw]{xcolor}
% Beamer presentation requires \usepackage{colortbl} instead of \usepackage[table,xcdraw]{xcolor}
\begin{table}[]
\begin{tabular}{|
>{\columncolor[HTML]{EFEFEF}}l |
>{\columncolor[HTML]{FFFFFF}}l |}
\hline
\textbf{22 septembre} & Introduction en XML et XML-TEI     \\ \hline
\textbf{29 septembre} & Structurer un document en TEI      \\ \hline
\textbf{06 octobre}   & XPath         \\ \hline
\textbf{13 octobre}   & XLST                               \\ \hline
\textbf{20 octobre}   & DST + Initiation JSON                    \\ \hline
\textbf{03 novembre}  & Exercices: XML et JSON avec Python \\ \hline
\end{tabular}
\end{table}
\end{frame}

\section{Introduction}
\begin{frame}{Récapitulatif des notions vues la semaine dernière}
\centering
    \textbf{\href{https://app.wooclap.com/events/JKAFURZ/questions/6aad0033ad4668ee6c1228cd}{Quiz sur Wooclap}}
\end{frame}

\begin{frame}{Qu'est-ce que le XML?}
\begin{block}{D'abord fut le GML... }
\vspace{1em}
    \textbf{\textit{Generalized Markup Language}}: développé par Charles Goldfarb, Edward Mosher et Raymond Lorie dans les années 1970.\\
    \textrightarrow{} Devient ensuite le Standard Generalized Markup Language (SGML) publié en 1986 et sert de base pour le HTML.\\
    \textrightarrow{} Problème: pas de normes pour utiliser le HTML.
    % Pourquoi ça a été crée?
\end{block}
\vspace{1em}
\begin{block}{... Puis apparut le XML}
\vspace{1em}
\textit{Extensible Markup Language}: publié en 1998 par le W3C (\textit{World Wide Web Consortium})\\ 
\textrightarrow{} Standard international.\\
\textrightarrow{} Sert de base à d'autres langages: XHTML, XLST etc...
\end{block}
    

\end{frame}


\begin{frame}{Pourquoi utiliser le XML?}

    \begin{itemize}
        \item Indiquer la fonction structurelle ou des éléments sémantiques = expliciter certains aspects du texte
        \item Format facile à lire par les machines et l'humain
        \item Interopérable
    \end{itemize}
% Sans ça, c'est nous qui identifions les informations importantes, pas l'ordinateur.
\end{frame}

\begin{frame}{Pourquoi utiliser le XML?}
Qu'est-ce qui peut être relevé dans un texte? \\
\pause 
\begin{itemize}
    \item Type de document: poèmes, essais, pièces de théâtre...
    \item Forme du document: paragraphes, vers, etc...
\end{itemize}
\textbf{Exemples:}
\begin{itemize}
    \item \href{https://lemdo.uvic.ca/qme\_editions/2.0/emdFV\_M.html}{\textit{Queen's Men Editions}}
    \item  \href{https://shelleygodwinarchive.org/sc/oxford/frankenstein/volume/i/\#/p1/mode/xml}{\textit{The Shelley-Godwin Archive}}
\end{itemize}
\end{frame}



\section{Structure d'un document XML}
\begin{frame}[fragile]{Encodage basique}
Le XML repose sur trois principes~:
\begin{itemize}
    \item Les balises sont appelées \textbf{éléments};
    \item Les éléments sont imbriqués;
    \item L'imbrication est sous forme d'arborescence
\end{itemize}
\vspace{1em}

\begin{tabular}{c}
    \begin{lstlisting}[numbers=none]
    <elementParent>
        <elementEnfant> Texte </elementEnfant>
    </elementParent>
    \end{lstlisting}
\end{tabular}

\end{frame}

\begin{frame}[fragile]{Encodage basique}
\begin{center}
    \begin{tabular}{c}
    \begin{lstlisting}[numbers=none]
    <Paragraphe>
        <Phrase> 
            <mot> Toto </mot>
        </Phrase>
    </Paragraphe>
    \end{lstlisting}
\end{tabular}
\end{center}
    
\end{frame}

\begin{frame}{Éléments et attributs}
\begin{block}{}
    \textbf{Élément =} Le contenu. Doit toujours être entouré de deux balises du même nom.\\
    Ex: \textcolor{blue}{<paragraphe></paragraphe>} \textrightarrow{} l'élément est le paragraphe.\\

    Les éléments ne doivent pas se chevaucher.\\
    Ex: \textcolor{blue}{<paragraphe>}\textcolor{BurntOrange}{<mot>}\textcolor{blue}{</paragraphe>}\textcolor{BurntOrange}{</mot>}
\end{block}

\begin{block}{}
    \textbf{Attribut =} Caractéristique du contenu. Ne peut pas réutilisé deux fois le même attribut.\\ 
    Ex: \textcolor{blue}{<paragraphe} \textcolor{red}{numero="1"}\textcolor{blue}{></paragraphe>} \textrightarrow{} c'est le premier paragraphe.
\end{block}
    
\end{frame}

\begin{frame}{Une structure bien formée}

Un encodage XML est dit \textcolor{red}{\textbf{bien formé}} s'il respecte les règles de balisage ci-dessous~:

\begin{itemize}
    \item Chaque balise doit avoir une balise de fin (une balise ouvrante et une balise fermante)~;
    \item Les éléments ne doivent pas se chevaucher et doivent bien s'imbriquer~;
    \item Un élément ne peut avoir qu'un seul même attribut~;
    \item Un seul élément racine (on doit avoir un seul élément qui englobe tous les autres).
\end{itemize}

\end{frame}


\begin{frame}[fragile]{Déclaration XML}
\begin{lstlisting}
        <?xml version="1.0" encoding="UTF-8"?>
    \end{lstlisting}

\begin{itemize}
    \item La déclaration XML s'écrit en début de document
    \item Elle indique la version utilisée de XML (toujours la version 1.0) et l'encodage du texte.
\end{itemize}
\end{frame}


\begin{frame}[fragile]{Autres commandes à connaître}
\begin{block}{Commenter le texte}
    \begin{lstlisting}
        <!-- on écrit son commentaire -->
    \end{lstlisting}
\end{block}

\begin{block}{Les entités}

\begin{verbatim}
    &entité;
\end{verbatim}
Permet d'utiliser des caractères interdits en XML comme le \&.
    
    
\end{block}
    
\end{frame}

\begin{frame}{Exercices}
    \url{https://pad.libreon.fr/eoi6u4UjRhO-XGzRv_bPzg?both}
\end{frame}



\section{Contraindre le XML}
\begin{frame}{Pourquoi contraindre le XML?}
    \begin{itemize}
        \item Uniformiser l'encodage~;
        \item Collaborer à plusieurs sur un même projet~;
        \item Automatisation.
    \end{itemize}
\end{frame}

\begin{frame}{Spécifications du XML: \textit{Document Type Definition} (DTD)}
    La \textbf{DTD} est une \textbf{liste de déclarations} qui permet de définir:
    \begin{itemize}
        \item la structure d'un document~;
        \item les éléments et attributs autorisés.
    \end{itemize}
    \textcolor{red}{\textrightarrow{} Un document XML peut être bien formé mais invalide car ne respecte par les règles de la DTD !}
\end{frame}

\begin{frame}[fragile]{Spécifications du XML: \textit{Document Type Definition} (DTD)}
    \begin{block}{Exemple (source W3C):}
\begin{lstlisting}
<!DOCTYPE NEWSPAPER [

<!ELEMENT NEWSPAPER (ARTICLE+)>
<!ELEMENT ARTICLE (HEADLINE,BYLINE,LEAD,BODY,NOTES)>
<!ELEMENT HEADLINE (#PCDATA)>
<!ELEMENT BYLINE (#PCDATA)>
<!ELEMENT LEAD (#PCDATA)>
<!ELEMENT BODY (#PCDATA)>
<!ELEMENT NOTES (#PCDATA)>

<!ATTLIST ARTICLE AUTHOR CDATA #REQUIRED>
<!ATTLIST ARTICLE EDITOR CDATA #IMPLIED>
<!ATTLIST ARTICLE DATE CDATA #IMPLIED>
<!ATTLIST ARTICLE EDITION CDATA #IMPLIED>
\end{lstlisting}

    \end{block}
        
\end{frame}

\begin{frame}{Spécifications du XML: \textit{Document Type Definition} (DTD)}
    Limites du DTD:
    \begin{itemize}
        \item N'est pas rédigé en XML~;
        \item Assez limité en terme d'attributs~;
        \item Difficile à lire~;
        \item Limité en espace de noms.
    \end{itemize}
\end{frame}

\begin{frame}{Les schémas XML: définition}
\begin{block}{}
\textbf{Définition:}\\
    \small{
    \og Un schéma spécifie \textbf{un ensemble de noms d'éléments}, \textbf{les noms et le type de données de tous les attributs} qui leur sont associés, et \textbf{les règles relatives aux contextes} dans lesquels ils peuvent apparaître.\fg{}}\\
    Burnard, Lou. \og La TEI et le XML \fg. \textit{Qu'est-ce que la Text Encoding Initiative }?, OpenEdition Press, 2015, \url{https://doi.org/10.4000/books.oep.1298}.
\end{block}
    
\end{frame}

\begin{frame}{Les schémas XML: exemple du XML-EAD}
    \textbf{\textit{Encoded Archival Description} =} crée dans les années 1990, c'est une DTD utilisée pour la description de fonds d'archive. Depuis 2007, est un format XML.\\
     \textbf{Exemple:} \href{https://ccfr.bnf.fr/portailccfr/servlet/ViewManager?menu=public_menu_view&record=eadbam:EADC:d0e1129770_FRBNFEAD0000129462&setCache=all_simple.PUBLIC_CATALOGUE_SIMPLE_MULTI&fromList=true&source=all_simple&action=public_search_federated&query=PUBLIC_CATALOGUE_SIMPLE}{Dossiers individuels de prisonniers de la Bastille}
\end{frame}




\section{XML-TEI}
\begin{frame}{Un peu d'histoire}
    \centering
    \small{
    \og La TEI a été d'abord développée, il y a plus de trente ans, comme un projet de recherche dans le champ alors émergent du \og Humanities computing \fg. L'idée originelle était de proposer un ensemble de recommandations sur la façon dont les chercheurs devraient créer des ressources textuelles \og lisibles par ordinateur\fg, qui soient adaptées aux besoins de la recherche - dans la mesure où un consensus existait sur le sujet -, mais qui soient également extensibles, puisque ces besoins changent et évoluent.\fg
    }\\
    \textbf{Burnard, Lou. \og Introduction \fg. \textit{Qu'est-ce que la Text Encoding Initiative ?}, OpenEdition Press, 2015, \url{https://doi.org/10.4000/books.oep.129}}
\end{frame}

\begin{frame}{Un peu d'histoire}
    \textbf{\textit{Text Encoding Initiative}:} créée en 1987 lors d'une conférence à Poughkeepsie sur les sciences humaines.\\
    Projet est de créer des recommandations sur l'encodage de textes informatiques.

    \begin{block}{Trois principes selon Lou Burnard:}
    \begin{itemize}
        \item le XML TEI s'intéresse au sens du texte plutôt qu'à son apparence;
        \item le XML TEI est indépendant de tout environnement logiciel particulier;
        \item le XML TEI a été conçu par la communauté scientifique, qui est aussi en charge de son développement continu.
    \end{itemize}
        
    \end{block}
\end{frame}

\begin{frame}{Tei Guidelines}
    TEI Consortium crée en 2000. Liste de recommandation: \url{https://tei-c.org/guidelines/}
\end{frame}

\begin{frame}[fragile]{TEI Guidelines}
    \begin{lstlisting}[numbers=none]
    <div>
        <p> 
            <persName> Toto </persName>
        </p>
    </div>
    \end{lstlisting}
\end{frame}

\begin{frame}[fragile]{Déclaration TEI}
    Pour utiliser les balises TEI, il faut déclarer l'espace de noms:\\
    \begin{lstlisting}
        <TEI xmlns="http://www.tei-c.org/ns/1.0">
    \end{lstlisting}
\end{frame}

\begin{frame}[fragile]{Squelette d'un document TEI}
    \begin{lstlisting}
        <TEI xmlns="http://www.tei-c.org/ns/1.0">
            <teiHeader>
            </teiHeader>
            <text>
                <body>
                    <p></p>
                </body>
            </text>
        </TEI>
    \end{lstlisting}
\end{frame}

\begin{frame}{Installations}
    Installer VS Code: \url{https://code.visualstudio.com/download?_exp_download=fb315fc982}\\
    Extension: Scholarly XML
\end{frame}

\begin{frame}[fragile]{Exercice: encoder en TEI}
Encoder à nouveau l'extrait du \textit{Bourgeois Gentilhomme} mais en XML-TEI en utilisant la documentation sur \href{https://www.tei-c.org/release/doc/tei-p5-doc/en/html/DR.html}{les textes théâtraux}\\
\begin{lstlisting}
    <castList> 
    <castItem> 
    <speaker> 
    <sp>
\end{lstlisting}    
\end{frame}


\begin{frame}{Exercice: encoder en TEI}
    Encoder le poème \textit{Heureux qui, comme Ulysse...} de Joachim Du Bellay en s'appuyant sur la documentation sur \href{https://www.tei-c.org/release/doc/tei-p5-doc/en/html/VE.html}{les poèmes}.
\end{frame}

\end{document}